\chapter{Troubleshooting}\label{app-troubleshooting}
This appendix collects problems that commonly arise before, during, or after a calculation. The input-script parser in \progname{} performs extensive structural and value checks. When an input is invalid, the message on standard error usually identifies the block or keyword involved and often states the allowed form. Correct the first reported input error before interpreting later messages, since one malformed block can cause several secondary diagnostics.

\section{The executable does not start}
If the executable terminates before an input script is read, first check the runtime environment. The binary release requires compatible Intel Fortran runtime libraries, Intel MKL, Intel MPI, Intel OpenMP, and standard Linux libraries. Use \texttt{ldd} on the executable to identify unresolved shared libraries. On a managed HPC system, load the site-provided Intel or oneAPI environment according to instructions provided by the administrator. On a workstation, confirm that the required library directories are visible through the environment described in Section~\ref{sec-install}. The installation script checks the environment but does not install third-party runtime libraries.

\section{An input block or keyword is rejected}
Check the number of dollar signs, the closing delimiters, and whether each nested block is placed under the correct parent job. Confirm that a keyword belongs to the selected job and that its value has the required data type. Integer sequences may use individual values, inclusive ranges such as \texttt{i..j} or \texttt{i:j}, stepped ranges such as \texttt{i:j:step}, and mixtures of these forms. A zero step is invalid.

Some older input names are no longer recognized. For orbital files, use \keyref{key-statepack-mofile}{mofile}. Select external electronic states with \keyref{key-statepack-state_index}{state\_index}. The parser reports context-specific diagnostics for malformed blocks, unknown keywords, invalid values, and missing required input. Read the first diagnostic on standard error and compare the input with the examples distributed with the program. Section~\ref{sec-stdout} shows how to redirect standard output and standard error to separate files when a permanent record is useful.

\section{Quasi-diabatic state character is ???}\label{sec-trouble-unphysical}
After a diabatization calculation, \progname{} uses the properties associated with the chosen method to characterize the resulting quasi-diabatic states. Where an assignment is available, standard output may identify a state as \texttt{LE at fragment 1}, \texttt{CT from fragment 2 to 3}, or another recognizable charge or excitation distribution. Occasionally the reported character is \texttt{???}. This indicates that the properties used for classification did not give an unambiguous physical assignment. In many cases, the intended quasi-diabatic states have not been constructed reliably, and the calculated couplings should not be used without further checks.

\paragraph{Check the selected adiabatic states.}
A diabatization requires adiabatic states that approximately span a sufficiently complete, reasonably self-contained Hamiltonian subspace for the physical process under study. If a relevant state has been omitted, a transformation of the selected states may be unable to recover the expected localized excitation or charge distribution. Inspect the energies, wave-function characters, and state indices printed during loading. Include missing states that are expected to mix appreciably with the selected states. Adding unrelated high-energy states is not a substitute for choosing a physically appropriate subspace.

\paragraph{Check whether the method matches the process.}
Each method distinguishes states through specific physical properties. GMH and FCD use dipole or fragment charge descriptors appropriate to charge-transfer problems. They are not suitable substitutes for an excitation-energy-transfer method when the relevant states have little charge-transfer character. FED is intended for two-fragment excitation-energy transfer, and the standalone FED job uses two target adiabatic states. Multistate FED--FCD addresses supported single-excitation problems involving both excitation and charge transfer. FPHD can distinguish a wider variety of particle and hole distributions, but the desired states must still be distinguishable by these quantities. Review the Background for the chosen method before interpreting an unexpected assignment.

\paragraph{Check fragments, the reference state, and the electronic-state data.}
An incorrect atom list may combine physically distinct fragments or split a chromophore unintentionally. Verify the atom numbering against the orbital file and inspect the fragment definitions. \textbf{Also confirm that the external calculation printed complete coefficients or amplitudes needed by the diabatization.} For FPHD and FED, check that the reference state is appropriate for the particle/hole or excitation-density comparison. If reference and target states come from separate calculations, confirm the molecular geometry, atom order, AO basis, and compatible wave-function representation. For FPHD, verify that the single- or multiple-excitation formulation is appropriate for the states loaded.

\paragraph{Examine the descriptors.}
An ambiguous assignment can occur when the calculated charge, excitation, or particle/hole descriptors do not clearly separate the anticipated state characters. Examine the numerical descriptors and their physical interpretation rather than relying only on the final character labels.
In an FPHD job, also inspect any reported Jacobi-optimization convergence information. The other diabatization methods do not generally use this optimization, so an unexplained assignment should not automatically be treated as an optimization-convergence problem.

%A reported character of \texttt{???} is a prompt to examine the inputs and the physical scope of the calculation. The absence of \texttt{???} does not by itself establish that the chosen adiabatic-state subspace or diabatization method is suitable. Interpret the resulting state characters alongside the relevant electronic structure and method-specific descriptors.

\section{An orbital or electronic-state file cannot be read}
Confirm that the path is correct from the working directory in which \progname{} was launched. Absolute and relative paths are both accepted. For a state package assembled from several calculations, check each package separately before diagnosing the combined calculation.

\progname{} identifies the external quantum chemistry program and the supported wave-function method from the data file. The \keyref{key-statepack-method}{method} keyword is normally unnecessary when only one supported method is present. It becomes important when the same log contains several supported data sets. For example, a Q-Chem output may contain both TDA and TDDFT results, while an ORCA output may contain CASSCF, NEVPT2, and QD-NEVPT2 information. If \keyref{key-statepack-method}{method} is omitted, the earliest supported method recognized in the file is used.

\section{The wrong electronic states were loaded}
The values in \keyref{key-statepack-state_index}{state\_index} refer to the order of states in the selected external method. They are not the internal state IDs used after loading. Section~\ref{sec-internal-ids} gives examples of this distinction. Check the state summary printed by \progname{} before using keywords such as \keyref{key-fphd-refid}{refid} or a \texttt{stateset} block.

For CIS, TDA, TDHF, TDDFT, and TDA-aTB data, external index 0 has a special meaning. It loads the underlying reference ground state as an independent state. Positive indices still refer to the explicitly stored states of the selected method. This makes an input such as \texttt{state\_index = 0..3} different from a four-state range in a method whose state numbering begins at 1.

\section{A custom Hamiltonian fails to load}
For \keyref{key-common-adiabatic_hamiltonian}{adiabatic\_hamiltonian}, use the numerical data format in Appendix~\ref{app-numfile}. The matrix dimension and element ordering must match the internal order of the adiabatic states in the corresponding diabatization. Every matrix element for this keyword is interpreted in eV.

The Phasefix \texttt{hamfile} keyword currently uses the separate indexed format in Appendix~\ref{app-hamfile}. A headerless Phasefix matrix must not be supplied as \texttt{adiabatic\_hamiltonian}, and a numerical data file with a \texttt{\#@} header must not be supplied to the current Phasefix reader.

\section{Cube output is too large or slow}
Cube generation can require substantial memory, disk space, and integration work. Request only the orbital or density cubes needed for interpretation. For preliminary visualization, a smaller \texttt{ngrid} can be useful. Increase the grid only when the required spatial resolution justifies the additional cost. Confirm that the selected output directory has sufficient free space before requesting many states or orbitals.
